\section{Literaturrecherche 3D-Netze}
\label{sec:literaturrecherche 3D-Netze}

\subsection{Einleitung}
Ein 3D-Netz bestimmt Lage und Höhe aller Netzpunkte in einer gemeinsamen Ausgleichung (Niemeier, 2008, S. 128, 247–248). 
GNSS liefert dabei 3D-Koordinaten oder Basislinien im geozentrischen System WGS84 (Niemeier, 2008, S. 330). 
Das Tachymeter (TPS) liefert Horizontalrichtungen, Zenitwinkel und Schrägdistanzen, die auf die lokale Lotrichtung bezogen sind (Banyai, 2011; França et al., 2021).
Die klassische Trennung in Lagenetz (2D) und Höhennetz (1D), wie sie etwa das Programm LTOP noch verwendet (swisstopo, o. J.), entfällt damit.
Der Bericht fasst zusammen, was die Literatur zu Aufbau, Modellierung, Ausgleichung und Qualitätsbeurteilung solcher Netze sagt. Hauptquellen sind das Lehrbuch von Niemeier (2008) zur Ausgleichungsrechnung und das Handbuch Ingenieurgeodäsie „Ingenieurbau“ (Möser et al., [2016]).
Ergänzt wird mit Fachartikeln und Praxisbeispielen aus der Schweiz.
Niemeier (2008, S. 128) nennt zwei Wege, die Messungen mit dem Koordinatensystem in Einklang zu bringen: Reduktion auf ein Abbildungssystem (z. B. LV95, UTM) oder Formulierung der Beobachtungsgleichungen direkt in einem räumlichen System. 
Für hochgenaue lokale Ingenieurnetze ist ein 3D-System von Bedeutung, wenn die Höhe streng mitverarbeitet werden soll. Für die Kombination von GNSS mit terrestrischen Messungen ist das globale kartesische System gebräuchlich.

\subsection{Messgrössen des 3D-Netzes}
\label{sec:Messgrössen des 3D-Netzes}
GNSS und TPS ergänzen sich folgendermassen: GNSS bestimmt Lagerung, Orientierung und Massstab des Netzes, Strecken- und Richtungsmessungen liefern die innere Geometrie mit hoher relativer Genauigkeit (Niemeier, 2008, S. 232, 247–248). 
Im Tunnelbau wird zum Beispiel ein GNSS-gestütztes Portalnetz mit Tachymetermessungen in den Berg übertragen (Leica Geosystems, 2016; Weber, 2014). 
Die folgende Tabelle zeigt die Messgrössen mit typischen Herstellerangaben aktueller Präzisionsinstrumente.

\begin{table}[h!]
    \centering
    \label{tab:messgrössen}
    \begin{tabularx}{\textwidth}{*{3}{>{\centering\arraybackslash}X}}
        \toprule
        \textbf{Messgrössen} & \textbf{Bezugssystem} & \textbf{Herstellergenauigkeit} \\
        \midrule
        GNSS-Basislinie $\Delta X$, $\Delta Y$, $\Delta Z$ (statisch) & geozentrisch (WGS84/ETRS89) & Lage 3 mm + 0,5 ppm, Höhe 5 mm + 0,5 ppm (Leica GS18) \\
        Horizontalrichtung (TPS) & lokal astronomisch (Lotrichtung) & $0{,}5^{\prime\prime} \approx 0{,}15\,\mathrm{mgon}$ (Leica TS60) \\
        Zenitwinkel (TPS) & lokal astronomisch (Lotrichtung) & $0{,}5^{\prime\prime} \approx 0{,}15\,\mathrm{mgon}$ (Leica TS60) \\
        Schrägdistanz (TPS) & systemunabhängig & 0,6 mm + 1 ppm (Leica TS60) \\
        Höhenunterschied (Nivellement) & Schwerefeld (Gebrauchshöhen) & Feldkurs Trimble Niv. NOCH Nachschauen \\
        \bottomrule
    \end{tabularx}
    \caption{Beispielgrössen in 3D-Netzen.}
\end{table}

Niemeier (2008, S. 330) betont, dass GNSS-Auswerteprogramme 3D-Koordinaten oder Basislinien im WGS84 liefern, die meist schon als Näherungswerte für die Netzausgleichung taugen.
Wichtig ist, die von der Software gelieferten voll besetzten Kofaktormatrizen in das stochastische Modell zu übernehmen, statt sie durch Diagonalmatrizen zu ersetzen (Niemeier, 2008, S. 249).
Die geringere Genauigkeit der GNSS-Höhenkomponente lässt sich so streng berücksichtigen (Niemeier, 2008, S. 365).
Die Herstellerangaben gelten für Idealbedingungen.
Im Netz kommen Zentrierfehler, Refraktion und Lotabweichung hinzu, die im funktionalen und stochastischen Modell berücksichtigt werden müssen (vgl. Abschnitt 3).

\subsection{Bezugssysteme und Modellbildung}
\label{sec:Bezugssysteme und Modellbildung}
GNSS und TPS messen in verschiedenen Systemen, aber ein 3D-Netz muss beide streng verbinden. 
Als gemeinsames System eignet sich das geozentrisch-kartesische System (X, Y, Z), in dem die terrestrischen Beobachtungsgleichungen nach den Formeln der «Dreidimensionalen Geodäsie» aufgestellt werden (Wolf, 1963, zitiert nach Niemeier, 2008, S. 247). 
Banyai (2011) zeigt diesen Ansatz für eine integrierte Ausgleichung von GNSS-Basislinien, Tachymeter- und Nivellementmessungen.
Ein Koordinatenunterschied im globalen System wird mit der Rotationsmatrix \(R(\Phi, \Lambda)\) in das lokale Horizontsystem \((n, e, u)\) gedreht. Daraus folgen Schrägdistanz und Zenitwinkel.
\[
\begin{pmatrix} n \\ e \\ u \end{pmatrix} = R(\Phi, \Lambda)\begin{pmatrix} \Delta X \\ \Delta Y \\ \Delta Z \end{pmatrix}, \qquad s = \sqrt{\Delta X^2 + \Delta Y^2 + \Delta Z^2}, \qquad z = \arccos\left(\frac{u}{s}\right)
\]

Bei der Lotabweichung wird die Abweichung zwischen der Lotrichtung und der Ellipsoide berücksichtigt. 
Das Tachymeter wird nach der Lotrichtung horizontiert, GNSS bezieht sich auf die Ellipsoide. Der Winkel zwischen beiden ist die Lotabweichung mit den Komponenten \(\xi\) (Nord-Süd) und \(\eta\) (Ost-West); astronomische und ellipsoidische Koordinaten hängen über \(\Phi = \varphi + \xi\) und \(\Lambda = \lambda + \eta / \cos\varphi\) zusammen.
Wird sie vernachlässigt, entsteht ein systematischer Fehler: In Brasilien ergaben sich bei einer Lotabweichung von $28{,}4^{\prime\prime}$ Lagefehler bis 45 ppm, also etwa 45 mm auf 1 km (França et al., 2021).
Banyai (2011) führt Lotabweichungen als Beobachtungen mit Standardabweichung $1^{\prime\prime}$ ein; bei Local-Tie-Netzen erreichen astrogeodätisch bestimmte Lotabweichungen etwa $0{,}2^{\prime\prime}$ (Suurmäki et al., 2022).
Im Gebirge, etwa in der Schweiz, ist dieser Effekt besonders relevant.\bigskip 

Die Krümmung des Zielstrahls verfälscht vor allem die Zenitwinkel und damit die trigonometrischen Höhen.
Üblich ist ein Refraktionskoeffizient k (Standardwert 0,13), der bei langen Visuren auch als Unbekannte geschätzt werden kann. 
Weber (2014) untersuchte im Portalnetz Sigirino des Ceneri-Basistunnels Refraktionseinflüsse auf lange Visuren, fand aber keine signifikante Verzerrung. \bigskip 

Bei den Höhen liefert GNSS die ellipsoidische Höhen h, die Praxis braucht Gebrauchshöhen. 
Die Verbindung läuft über die Geoidundulation N mit $H = h - N$; in der Schweiz über das Geoidmodell CHGeo (swisstopo, 2004).
Viele Programme der Praxis, z. B. LTOP von swisstopo, rechnen Lage und Höhe getrennt (2D+1D) und sind damit keine strengen 3D-Ausgleichungen (swisstopo, o. J.).

\subsection{Ausgleichung im Gauss-Markov-Modell}
\label{sec:Ausgleichung im Gauss-Markov-Modell}

Ausgeglichen wird nach der Methode der kleinsten Quadrate im Gauss-Markov-Modell: Unbekannte sind die 3D-Koordinaten aller Punkte plus Zusatzparameter (Orientierungen, Massstab, Refraktion, Transformationsparameter). 
Das funktionale Modell verknüpft Beobachtungen und Unbekannte, das stochastische Modell beschreibt deren Genauigkeit (Niemeier, 2008, Kap. 4):
\[
\begin{aligned}
\mathbf{l} + \mathbf{v} &= \mathbf{A}\,\hat{\mathbf{x}}, &
\boldsymbol{\Sigma}_{ll} &= \sigma_0^2\,\mathbf{Q}_{ll}, \\
\hat{\mathbf{x}} &= \left(\mathbf{A}^{T}\mathbf{P}\mathbf{A}\right)^{-1}
\                    \mathbf{A}^{T}\mathbf{P}\,\mathbf{l}, &
\mathbf{P} &= \mathbf{Q}_{ll}^{-1}
\end{aligned}
\]

Die Näherungskoordinaten: Die Beobachtungsgleichungen sind nichtlinear und werden um Näherungswerte linearisiert. 
Für terrestrische 3D-Netze ist eine direkte 3D-Näherungsberechnung numerisch instabil. Stabil ist die Zerlegung in Lage und Höhe, danach folgt die gemeinsame 3D-Ausgleichung mit Raumstrecken und Zenitwinkeln. 
GNSS-Koordinaten sind dagegen meist direkt brauchbare Näherungswerte (Niemeier, 2008, S. 330 f.).\bigskip

Der Datumsdefekt: Messungen legen nur die innere Geometrie fest, nicht die Lagerung im Koordinatensystem. 
Ein 3D-Netz hat maximal 7 Datumsdefekte: 3 Translationen, 3 Rotationen und den Massstab (Niemeier, 2008, S. 232, Tab. 7.1). 
Strecken liefern den Massstab, GNSS zusätzlich Orientierung und Lage.
Das datumsfreie Konzept: Für die strenge gemeinsame Ausgleichung von GNSS und TPS empfiehlt Niemeier (2008, S. 247–249), je Beobachtungsgruppe eigene Zusatzparameter einzuführen, nämlich je GNSS-Session 3 Translationen, je Gruppe von Zenitwinkeln, Azimuten und GNSS-Daten eigene Rotationen und je Instrument einen Massstab. 
Damit läuft die Ausgleichung mit dem maximalen Defekt d = 7, und die innere Netzqualität ist klar von der Netzanbindung getrennt. 
Soll GNSS oder ein EDM den verbindlichen Netzmassstab liefern, lässt man für diese Gruppe den Massstabsparameter weg. \bigskip 

Die Datumsfestlegung: Damit wird das Datum festgelegt (Niemeier, 2008, Kap. 7.3):
1. Zwangsfreie Lagerung: genau d Koordinaten werden festgehalten; die innere Geometrie bleibt unverändert.
2. Gesamt- oder Teilspurminimierung (freies Netz): das Netz wird auf alle oder ausgewählte Punkte gelagert; liefert die kleinste Spur der Kovarianzmatrix.
3. Hierarchische Ausgleichung und weiche Lagerung: Anschlusspunkte werden fest bzw. als Beobachtungen mit Genauigkeit eingeführt. Die innere Geometrie wird dabei verändert. \bigskip 

Für die Integration von GNSS in ein bestehendes Netz beschreibt Niemeier (2008, S. 363–365) ein zweistufiges Vorgehen:
Zuerst Näherungstransformation mit bekannten Parametern, dann gemeinsame Ausgleichung mit einer differenziellen 7-Parameter-Helmert-Transformation.
Pro GNSS-Datensatz ein eigener Parametersatz.
Ein Praxisbeispiel ist das Ceneri-Portalnetz. Dort zeigte sich, dass durch die GNSS-Transformation ins LV03 ein Massstabsfaktor von $-14\,\mathrm{ppm}$ entstand, der bestätigt wurde; ausgeglichen wurde mit weicher Lagerung in der Lage und gezwängter Lagerung in der Höhe (Weber, 2014). \bigskip 

Das stochastisches Modell und Varianzkomponenten:
Jede Beobachtungsgruppe (Richtungen, Strecken, Zenitwinkel, GNSS) erhält einen eigenen Varianzfaktor, der mit einer Varianzkomponentenschätzung bestimmt wird (Niemeier, 2008, S. 249, Kap. 9.3). 
So wird die Gewichtung der heterogenen Messgrössen aus den Daten selbst kalibriert. 
Wie heikel eine falsche Gewichtung ist, zeigt Niemeiers Beispiel Vianden: Eine einzelne Massstabsstrecke am Netzrand vergrösserte die Konfidenzellipsen deutlich, auch bei zehnfach höherer Genauigkeit (Niemeier, 2008, S. 245–247).

\subsection{Genauigkeit und Zuverlässigkeit}
\label{sec:Genauigkeit und Zuverlässigkeit}
Die Qualität eines 3D-Netzes wird an zwei Kriterien gemessen. Erstens an der Genauigkeit, also daran, wie stark die Koordinaten streuen, und an der Zuverlässigkeit, also daran, wie gut grobe Fehler aufgedeckt werden (Niemeier, 2008, Kap. 8). 
Am Anfang steht der Globaltest, der mit einem $\chi^2$- bzw. F-Test prüft, ob die a-posteriori-Varianz $\hat{\sigma}_0^2$ mit der a-priori-Varianz $\sigma_0^2$ verträglich ist. Wird er abgelehnt, deutet dies auf grobe Fehler oder auf ein unzutreffendes funktionales oder stochastisches Modell hin (Niemeier, 2008, S. 167 ff.). 
Die Genauigkeit wird im 3D-Netz nicht mehr mit Fehlerellipsen, sondern mit Konfidenzellipsoiden je Punkt beschrieben, die sich aus der $3 \times 3$-Teilmatrix von $\Sigma_{xx}$ ergeben. 
Als skalare Masse dienen der Helmertsche Punktfehler, der auf der Spur der Kovarianzmatrix beruht, und der Werkmeistersche Punktfehler, der im 3D-Fall dem Volumen des Fehlerellipsoids entspricht. Alle diese Massen hängen allerdings von der Datumsfestlegung ab (Niemeier, 2008, S. 274--286). 
Für die innere Zuverlässigkeit ist der Redundanzanteil $r_i$ zentral, der angibt, welchen Anteil eine Beobachtung an der Gesamtredundanz trägt (Förstner, 1979, zitiert nach Niemeier, 2008, S. 293). 
Im Mittel beträgt er $1-u/n$, bei 30 Beobachtungen für 20 Unbekannte also nur 0,33, und Werte nahe 0 zeigen, dass ein Fehler in dieser Beobachtung kaum sichtbar wird (Niemeier, 2008, S. 293). 
Nach Baarda (1968, zitiert nach Niemeier, 2008, S. 295) wird jede Beobachtung anhand ihrer normierten Verbesserung getestet. Dies wird auch Data Snooping genannt. \bigskip 

Dies ist ein Verfahren nach Baarda (1968, zitiert nach Niemeier, 2008, S. 294--295), mit dem nach der Ausgleichung jede einzelne Beobachtung auf grobe Fehler geprüft wird. 
Dazu wird für jede Beobachtung $l_i$ die standardisierte Verbesserung $w_i = v_i / \sigma_{v_i}$ berechnet. Ist das Ausgleichungsmodell korrekt, folgt sie der Standardnormalverteilung $N(0,1)$ (Niemeier, 2008, S. 293--294). 
Ein Ausreisser liegt vor, wenn $\lvert w_i \rvert$ einen aus der Normalverteilung abgeleiteten Grenzwert übersteigt. Bei einer Irrtumswahrscheinlichkeit von $\alpha_0 = 0{,}1\,\%$ liegt dieser bei rund 3,29 (Niemeier, 2008, S. 294). 
Grundlage ist die Annahme, dass im gesamten Ausgleichungsproblem nur ein einzelner Ausreisser vorliegt (Baarda, 1968, zitiert nach Niemeier, 2008, S. 294). 
Weil sich ein grober Fehler über sogenannte Verschmierungseffekte auch auf die Verbesserungen benachbarter Beobachtungen auswirkt, geht man iterativ vor: 
Es wird jeweils nur die Beobachtung mit dem grössten $w_i$ als grob falsch eingestuft, danach wird neu ausgeglichen und erneut getestet (Niemeier, 2008, S. 222). 
Besonders problematisch sind Beobachtungen mit kleinem Redundanzanteil $r_i$, denn ein Fehler in ihnen kann sich stärker auf andere Verbesserungen auswirken als auf die eigene. 
Niemeier empfiehlt deshalb, bei der Netzplanung Redundanzanteile unter 0,2 zu vermeiden (Niemeier, 2008, S. 295).\bigskip 

Wie gross ein Fehler mindestens sein muss, damit er mit einer Testgüte von $\gamma_0 = 80\,\%$ aufgedeckt wird, gibt der minimal aufdeckbare Fehler $\nabla l_i$ an (Niemeier, 2008, S. 295--296).
Mit einer Irrtumswahrscheinlichkeit von $\alpha_0 = 0{,}1\,\%$ und einer Testgüte von $\gamma_0 = 80\,\%$ ergibt sich der Nichtzentralitätsparameter $\delta_0 = 4{,}13$ und daraus der minimal aufdeckbare Fehler $\nabla l_i = \delta_0 \cdot \sigma_{l_i} / \sqrt{r_i}$. 
Bei GNSS-Basislinien wird dabei nicht eine einzelne Komponente, sondern der ganze 3D-Vektor als fehlerhaft getestet, was einer dreidimensionalen Alternativhypothese entspricht (Niemeier, 2008, Kap. 8.3). 
Die äussere Zuverlässigkeit schliesslich beschreibt, wie stark ein gerade nicht mehr aufdeckbarer Fehler die Koordinaten verfälscht, nämlich um $\nabla x = Q_{xx} A^T P \nabla l$. 
Für die Nutzerinnen und Nutzer eines Netzes ist diese maximale Koordinatenverfälschung oft aussagekräftiger als der Redundanzanteil selbst, sie ist jedoch ebenfalls datumsabhängig (Niemeier, 2008, S. 304 f.). 
Da alle genannten Masse nur von der Netzkonfiguration und den angenommenen Genauigkeiten abhängen, lassen sie sich bereits in der Planung durch Simulation berechnen und zur Optimierung des Netzes nutzen (Niemeier, 2008, Kap. 9.5).

\subsection{Fazit zur Literaturrecherche}
\label{sec:Fazit zur Literaturrecherche}

Die Literaturrecherche zeigt, dass die theoretischen Grundlagen von 3D-Netzen gut dokumentiert sind. 
Das Lehrbuch von Niemeier (2008) deckt Modellbildung, Datumsfestlegung, Qualitätsbeurteilung und die gemeinsame Ausgleichung von GNSS- und terrestrischen Messungen zusammenhängend ab und ist deshalb die Hauptquelle dieses Berichts.
Das Handbuch von Möser liefert den Rahmen für die Anwendung im Ingenieurbau.
Neuere Entwicklungen wie aktuelle Instrumente oder frei verfügbare 3D-Software sind in diesen Lehrbüchern noch nicht enthalten.
Diese Lücke schliessen jüngere Fachartikel und Konferenzbeiträge, etwa zur integrierten 3D-Ausgleichung (Banyai, 2011), zum Einfluss der Lotabweichung (França et al., 2021) und zur Datumsfestlegung in Local-Tie-Netzen (Suurmäki et al., 2022). 
Die Praxisbeispiele aus der Schweiz stammen vor allem aus Bachelorarbeiten (Ambrosini \& Merz, 2022; Weber, 2014) und einer Firmenpublikation (Leica Geosystems, 2016).
Sie sind anschaulich, aber nicht wissenschaftlich begutachtet und dienen deshalb nur als Beispiele.
Die Genauigkeitsangaben beruhen auf Herstellerdatenblättern (Leica Geosystems, o. J., 2025) und sind als Idealwerte zu verstehen. 
Auffällig ist ein Unterschied zwischen Theorie und Praxis, weil in der Literatur klar die strenge 3D-Ausgleichung empfohlen wird, die Schweizer Praxis rechnet mit LTOP aber häufig noch getrennt nach Lage und Höhe (swisstopo, o. J.). 
Begutachtete Arbeiten zu kombinierten GNSS-TPS-3D-Netzen in der Schweiz wurden in dieser Recherche kaum gefunden.
Für diesen Unterschied gibt es mehrere Gründe. 
Der wichtigste dürfte das Höhensystem sein weil offiziell gilt in der Schweiz weiterhin LN02, ein Gebrauchshöhensystem ohne Schwerekorrekturen, das sich nicht streng mit ellipsoidischen GNSS-Höhen verbinden lässt (swisstopo, 2004; Willi et al., 2024). 
Ein neues Höhensystem ist erst ab 2028 geplant (Willi et al., 2024). 
Zudem werden Höhen mit dem Präzisionsnivellement genauer bestimmt als mit GNSS (GeoConnexion, 2016; Leica Geosystems, 2025), und LTOP liefert den für die amtliche Vermessung geforderten Genauigkeits- und Zuverlässigkeitsnachweis (swisstopo, o. J()). 
Eine 3D-Ausgleichung bedeutet ausserdem mehr Aufwand (Niemeier, 2008, S. 330), während der Einfluss der Lotabweichung erst bei grossen Lotabweichungen und steilen Visuren bedeutend wird (França et al., 2021). 
Dass wenig dazu publiziert ist, heisst jedoch nicht, dass 3D nicht angewendet wird: Streng dreidimensionale Programme wie Trinet+ sind in der Praxis im Einsatz (Vasi \& Meyenberg, 2014), die Ergebnisse stehen aber meist in internen Projektberichten.